\documentclass[11pt]{article}

\usepackage[a4paper,margin=1in]{geometry}
\usepackage{amsmath,amssymb,amsthm,mathtools,bm}
\usepackage[T1]{fontenc}
\usepackage{lmodern}
\usepackage{microtype}
\usepackage{enumitem}
\usepackage{booktabs}
\usepackage{hyperref}
\usepackage{fancyhdr}

\hypersetup{
  colorlinks=false,
  hidelinks
}

\setlength{\parindent}{0pt}
\setlength{\parskip}{0.55em}
\setlist[itemize]{leftmargin=1.5em}
\setlist[enumerate]{leftmargin=1.7em}

\pagestyle{fancy}
\fancyhf{}
\lhead{\textsc{Lecture 1}}
\rhead{\thepage}
\renewcommand{\headrulewidth}{0.4pt}

\newcommand{\ket}[1]{\lvert #1\rangle}
\newcommand{\bra}[1]{\langle #1\rvert}
\newcommand{\braket}[2]{\langle #1\mid #2\rangle}
\newcommand{\Prob}{\operatorname{Pr}}
\newcommand{\E}{\mathbb{E}}
\newcommand{\Var}{\operatorname{Var}}
\newcommand{\C}{\mathbb{C}}
\newcommand{\R}{\mathbb{R}}

\begin{document}

\begin{center}
  {\LARGE\bfseries Lecture 1}\\[0.35em]
  {\large September 24, 2025}
\end{center}

\hrule
\vspace{0.8em}

\section{The Two Faces of Quantum Computation}

There are two complementary aspects of quantum computation:

\begin{enumerate}
  \item \textbf{The theoretical model of quantum computation.}
  \begin{itemize}
    \item An idealized, consistent, plausible description of reality.
    \item ``Implementation agnostic.''
    \item A testbed for algorithm design and proving impossibility results.
  \end{itemize}

  \item \textbf{Implementations of quantum computation.}
  \begin{itemize}
    \item Superconducting qubits, photonic systems, neutral atoms, etc.
    \item Models in which we hope to simulate the idealized model.
    \item Not the main focus of this class, but pertinent to the implementation
          and manipulation of quantum information.
  \end{itemize}
\end{enumerate}

An example of quantum-computing advantage.

\section{Notation}

We express vectors over $\C^d$ as
\[
  \ket{v}
  =
  \begin{pmatrix}
    v_0\\
    v_1\\
    \vdots\\
    v_{d-1}
  \end{pmatrix}.
\]

Because of its prevalence, we write
\[
  \ket{v}
  =
  v_0\ket{0}+v_1\ket{1}+\cdots+v_{d-1}\ket{d-1}.
\]

The complex conjugate is denoted by $(\cdot)^*$.

For vectors,
\[
  \bra{v}
  =
  \begin{pmatrix}
    v_0^* & v_1^* & \cdots & v_{d-1}^*
  \end{pmatrix},
\qquad
  \braket{v}{v}
  =
  \sum_{j=0}^{d-1}|v_j|^2.
\]

\textbf{Questions for class.}
\begin{enumerate}
  \item What is $\braket{1}{1}$?
  \item Let $M=\ket{1}\bra{1}$ (short form $|1\rangle\langle1|$).
        For unit vectors $\ket{v}$, what is $M\ket{v}$?
\end{enumerate}

\section{Axiom 1: State Description}

The state of exactly one qubit can be expressed as a vector in $\C^2$
of unit norm:
\[
  \ket{\psi}
  =
  \begin{pmatrix}
    \alpha\\
    \beta
  \end{pmatrix},
  \qquad
  |\alpha|^2+|\beta|^2=1.
\]

Define
\[
  \ket{0}=
  \begin{pmatrix}
    1\\0
  \end{pmatrix},
  \qquad
  \ket{1}=
  \begin{pmatrix}
    0\\1
  \end{pmatrix}.
\]

Note that
\[
  \braket{0}{1}=0.
\]

These form an orthonormal basis for $\C^2$ and are possible states for a qubit.

If the state of a qubit is either $\ket{0}$ or $\ket{1}$, we call it
\emph{classical}. For example,
\[
  \ket{0}+\ket{1}
  =
  \begin{pmatrix}
    1\\1
  \end{pmatrix}
\]
(up to normalization) is a non-classical, i.e.\ ``quantum,'' possible state.

\section{Axiom 2: Transformation of a Qubit}

A quantum transformation is represented by a unitary matrix $U$:
\[
  \ket{\psi}\longmapsto U\ket{\psi}.
\]

For a unitary matrix,
\[
  U^\dagger U=I,
\]
and consequently a unit vector remains a unit vector:
\[
  \|U\ket{\psi}\|=\|\ket{\psi}\|.
\]

\section{Single-Qubit Gates}

Important examples include:

\[
  X=
  \begin{pmatrix}
    0&1\\
    1&0
  \end{pmatrix}
  \qquad\text{``bit flip,''}
\]

\[
  Z=
  \begin{pmatrix}
    1&0\\
    0&-1
  \end{pmatrix}
  \qquad\text{``phase flip,''}
\]

and the Hadamard gate
\[
  H=
  \frac{1}{\sqrt2}
  \begin{pmatrix}
    1&1\\
    1&-1
  \end{pmatrix}
  \qquad\text{``Hadamard.''}
\]

\section{Axiom 3: Measurement / Born's Rule}

Given a quantum state
\[
  \ket{\psi}
  =
  \alpha\ket{0}+\beta\ket{1}
  =
  \begin{pmatrix}
    \alpha\\
    \beta
  \end{pmatrix},
\]
we can measure the quantum state.

The probability of observing $0$ is
\[
  \Prob(0)=|\alpha|^2,
\]
in which case the state collapses to $\ket{0}$.

The probability of observing $1$ is
\[
  \Prob(1)=|\beta|^2,
\]
in which case the state collapses to $\ket{1}$.

The classical value ``0'' or ``1'' is output.

\subsection{Repeated Measurement}

If the state is $\ket{0}$ or $\ket{1}$, measurement does not change the state.
Hence $\ket{0}$ and $\ket{1}$ behave as classical values: measurement /
observation does not change the state.

\section{Global Phase}

Consider the states
\[
  \ket{\psi}
  \quad\text{and}\quad
  e^{i\theta}\ket{\psi}.
\]

These states cannot be distinguished by measurement. For example,
\[
  \Prob\!\left[
    e^{i\theta}\ket{\psi}\text{ collapsing to }\ket{j}
  \right]
  =
  \left|\bra{j}e^{i\theta}\ket{\psi}\right|^2
  =
  \left|e^{i\theta}\bra{j}\ket{\psi}\right|^2
  =
  |\bra{j}\ket{\psi}|^2.
\]

Thus $e^{i\theta}$ is called the \emph{global phase}, and quantum states
form an equivalence class
\[
  \ket{\psi}\sim e^{i\theta}\ket{\psi}.
\]

The set of physical pure states is therefore the corresponding quotient
space (projective Hilbert space).

\section{Axiom 4: Initialization}

We can initialize a qubit as
\[
  \ket{\psi_0}=\ket{0}.
\]

These axioms already imply a perfect random-number generator.

\begin{enumerate}
  \item Initialize the qubit as $\ket{\psi_0}=\ket{0}$.
  \item Apply the Hadamard transform:
  \[
    \ket{\psi_1}
    =
    H\ket{0}
    =
    \frac{1}{\sqrt2}\bigl(\ket{0}+\ket{1}\bigr).
  \]
  \item Measure the qubit:
  \[
    \Prob(0)=\frac12,
    \qquad
    \Prob(1)=\frac12.
  \]
\end{enumerate}

\section{Multiple Qubits}

Recall the notion of tensor products.

For matrices $A$ and $B$,
\[
  A\otimes B
\]
is the Kronecker (tensor) product. If
$A=(a_{ij})$ and $B=(b_{kl})$, then
\[
  A\otimes B
  =
  \begin{pmatrix}
    a_{11}B & \cdots & a_{1n}B\\
    \vdots  & \ddots & \vdots\\
    a_{m1}B & \cdots & a_{mn}B
  \end{pmatrix}.
\]

\subsection{Properties of the Tensor Product}

For compatible matrices,
\begin{enumerate}
  \item Associativity:
  \[
    (A\otimes B)\otimes C
    =
    A\otimes(B\otimes C).
  \]
  \item Linearity:
  \[
    (A\otimes B)(C\otimes D)
    =
    (AC)\otimes(BD).
  \]
  \item Transpose:
  \[
    (A\otimes B)^T=A^T\otimes B^T.
  \]
  \item Conjugate transpose:
  \[
    (A\otimes B)^\dagger
    =
    A^\dagger\otimes B^\dagger.
  \]
  \item When invertible,
  \[
    (A\otimes B)^{-1}
    =
    A^{-1}\otimes B^{-1}.
  \]
\end{enumerate}

We can also take tensor products of vectors, since vectors are $1$-D matrices.

\subsection{Exercises}

\begin{enumerate}
  \item What is $\ket{0}\otimes\ket{1}$?
  \item What is $\ket{0}\otimes\ket{0}\otimes\ket{1}$?
\end{enumerate}

\section{Computational-Basis Notation}

We use the shorthand
\[
  \ket{0}\ket{1}\ket{0}
  =
  \ket{010}.
\]

More generally,
\[
  \ket{x_1x_2\cdots x_n},
  \qquad x_i\in\{0,1\},
\]
denotes the computational-basis state whose bit string is
$x_1x_2\cdots x_n$.

For $j\in\{0,\ldots,2^n-1\}$, we write
\[
  \ket{j}
\]
for the basis vector corresponding to the binary representation of $j$.

Equivalently,
\[
  \ket{j}
  =
  \ket{j_1j_2\cdots j_n},
\]
where $j_1j_2\cdots j_n$ is the $n$-bit binary representation of $j$.
These basis states are orthogonal:
\[
  \braket{j}{k}=\delta_{jk}.
\]

For $\ket{j}\in\C^d$,
\[
  \ket{j}
  =
  \begin{pmatrix}
    0\\
    \vdots\\
    1\\
    \vdots\\
    0
  \end{pmatrix},
\]
with the $1$ in the $j$th position (up to the indexing convention).

\section{Bringing Multiple Qubits Together}

Suppose qubit $A$ is in state $\ket{\psi_A}$ and qubit $B$ is in state
$\ket{\psi_B}$. We describe the state of the two qubits by
\[
  \ket{\psi_{AB}}
  =
  \ket{\psi_A}\otimes\ket{\psi_B}.
\]

We need axioms to describe the composite system.

\begin{enumerate}
  \item If qubits do not interact, the description should reproduce the
        statistics of the individual qubits.
  \item There must be a method for entangling the qubits; i.e.\ for any
        desired state there should exist a method of generating that state.
\end{enumerate}

\section{Axioms for $n$ Qubits}

\begin{enumerate}
  \item The state of an $n$-qubit system is a unit vector in
  \[
    (\C^2)^{\otimes n}
    \cong \C^{2^n},
  \]
  the $n$-fold tensor-product Hilbert space.

  Every quantum state is therefore of the form
  \[
    \ket{\psi}
    =
    \sum_{x\in\{0,1\}^n}\alpha_x\ket{x},
    \qquad
    \sum_x|\alpha_x|^2=1.
  \]

  \item For any unitary $U\in\C^{2^n\times2^n}$, we can transform
  \[
    \ket{\psi}\longmapsto U\ket{\psi}.
  \]

  In particular, products of local and two-qubit gates can generate
  general quantum transformations. The lecture notes emphasize that
  allowing gates acting on adjacent qubits is sufficient for the intended
  construction.

\end{enumerate}

\subsection{Classical States and Superposition}

A state
\[
  \ket{\psi}
  =
  \sum_x\alpha_x\ket{x}
\]
is \emph{classical} if all amplitudes are zero except one (whose magnitude
is $1$).

A state is in \emph{superposition} if there are multiple non-zero
amplitudes.

\section{Local Gates and Two-Qubit Gates}

A unitary in $\C^{2\times2}$ is a single-qubit unitary. In particular,
$H\in\C^{2\times2}$ is a single-qubit unitary.

A unitary in $\C^{4\times4}$ is a two-qubit unitary.

A circuit of the form
\[
  I\otimes\cdots\otimes I\otimes H\otimes I\otimes\cdots\otimes I
\]
acts on one qubit, while a two-qubit gate acts on a pair of qubits.

More generally, a product of such local gates and two-qubit gates generates
a unitary transformation on the full $n$-qubit space.

\subsection{Linearity}

Suppose
\[
  U
  =
  I\otimes\cdots\otimes I\otimes H\otimes I\otimes\cdots\otimes I.
\]

For a computational-basis state
\[
  \ket{x_1x_2\cdots x_n},
\]
the action factors according to the qubit on which $H$ acts. By linearity,
for
\[
  \ket{\psi}
  =
  \sum_{x_1,\ldots,x_n}
  \alpha_{x_1\cdots x_n}
  \ket{x_1\cdots x_n},
\]
we obtain
\[
  U\ket{\psi}
  =
  \sum_{x_1,\ldots,x_n}
  \alpha_{x_1\cdots x_n}
  U\ket{x_1\cdots x_n}.
\]

\section{Measurement of All Qubits}

The $n$-qubit version of Born's rule is:

\[
  \ket{\psi}
  =
  \sum_{j=0}^{2^n-1}\alpha_j\ket{j}
  \quad\Longrightarrow\quad
  \Prob(j)=|\alpha_j|^2.
\]

Upon observing outcome $j$, the state collapses to
\[
  \ket{j}.
\]

The problem set will introduce how to measure arbitrary single qubits.

\subsection{Measurement of the First Qubit}

For an $n$-qubit state, one can write
\[
  \ket{\psi}
  =
  \ket{0}\ket{\psi_0}
  +
  \ket{1}\ket{\psi_1}.
\]

Measuring the first qubit gives the corresponding classical outcome and
leaves the remaining system in the appropriate conditional state.

\subsection{Adding a Fresh Qubit}

Given an $n$-qubit state $\ket{\psi}$, we can initialize a fresh qubit as
$\ket{0}$, generating
\[
  \ket{\psi}\otimes\ket{0},
\]
a state on $n+1$ qubits.

\section{Joining and Separating Systems}

More generally, we can join quantum systems:
\[
  \ket{\psi}\otimes\ket{\phi}.
\]

We can also separate unentangled systems. These are not additional axioms;
they can be derived from the axioms.

\subsection{Why the Description Gives the Correct Statistics}

Suppose
\[
  \ket{\psi}
  =
  \alpha\ket{0}+\beta\ket{1},
\]
where
\[
  |\alpha|^2+|\beta|^2=1.
\]

Then
\[
  \ket{\psi}\otimes\ket{\phi}
  =
  \alpha\ket{0}\otimes\ket{\phi}
  +
  \beta\ket{1}\otimes\ket{\phi}.
\]

If we measure the first qubit and obtain $1$, the remaining state is
proportional to $\ket{\phi}$, as expected for an unentangled system.

\section{Important Math Review: Probability Theory}

Let $X\ge 0$ be a random variable.

\subsection{Markov's Inequality}

For all $a>0$,
\[
  \Prob(X\ge a)
  \le
  \frac{\E[X]}{a}.
\]

Also,
\[
  \E[X]=\sum_x x\,\Prob(X=x)
\]
for a discrete random variable.

\subsection{Chebyshev's Inequality}

Let $X$ be a random variable with finite variance
\[
  \Var(X)=\sigma^2
\]
and mean
\[
  \E[X]=\mu.
\]

For $k>0$,
\[
  \Prob\!\left(|X-\mu|\ge k\sigma\right)
  \le
  \frac{1}{k^2}.
\]

This follows by applying Markov's inequality to
\[
  Y=(X-\mu)^2
\]
with threshold $a=k^2\sigma^2$.

\subsection{Chernoff Bounds}

For a random variable $X$ and suitable $t>0$,
\[
  \Prob(X\ge a)
  =
  \Prob(e^{tX}\ge e^{ta})
  \le
  e^{-ta}\E[e^{tX}],
\]
by Markov's inequality.

If
\[
  X=X_1+\cdots+X_n
\]
for independent random variables, then
\[
  \E[e^{tX}]
  =
  \E\!\left[e^{t\sum_iX_i}\right]
  =
  \prod_i \E[e^{tX_i}].
\]

\subsection{Bernoulli Case}

Let $X_i\in\{0,1\}$ with
\[
  \Prob(X_i=1)=p_i,
  \qquad
  \Prob(X_i=0)=1-p_i.
\]

Then
\[
  \E[e^{tX_i}]
  =
  (1-p_i)+p_i e^t
  =
  1+p_i(e^t-1).
\]

Using
\[
  1+x\le e^x,
\]
we obtain
\[
  \E[e^{tX_i}]
  \le
  \exp\!\bigl(p_i(e^t-1)\bigr).
\]

Therefore,
\[
  \E[e^{tX}]
  \le
  \exp\!\left((e^t-1)\sum_i p_i\right).
\]

Since
\[
  \mu=\E[X]=\sum_i p_i,
\]
we get
\[
  \E[e^{tX}]
  \le
  e^{\mu(e^t-1)}.
\]

Consequently,
\[
  \Prob(X\ge a)
  \le
  \inf_{t>0}
  \exp\!\left(-ta+\mu(e^t-1)\right).
\]

Taking
\[
  a=(1+\delta)\mu
\]
and optimizing gives
\[
  t=\ln(1+\delta).
\]

This yields the standard Chernoff upper-tail bound
\[
  \Prob\!\left(X\ge(1+\delta)\mu\right)
  \le
  \left(
    \frac{e^\delta}{(1+\delta)^{1+\delta}}
  \right)^\mu.
\]

A common simpler form is
\[
  \Prob\!\left(X\ge(1+\delta)\mu\right)
  \le
  \exp\!\left(-\frac{\delta^2\mu}{2+\delta}\right),
  \qquad \delta>0.
\]

Similarly, for $0<\delta<1$,
\[
  \Prob\!\left(X\le(1-\delta)\mu\right)
  \le
  \exp\!\left(-\frac{\delta^2\mu}{2}\right).
\]

\subsection*{Exercise}

Let
\[
  X=\sum_{i=1}^n X_i,
  \qquad
  \mu=\E[X].
\]
Derive the corresponding Chernoff bounds directly from Markov's inequality.

\end{document}
